\section{模型评价与推广}

\subsection{模型优点}

本文建立的圆柱镜反射逆映射模型在理论完备性和实用性方面均具有显著优势。

逆映射公式具有解析闭式解，计算效率高，无需数值迭代即可直接求得纸面坐标。对于分辨率为 $N_\theta \times N_h$ 的圆柱展开面，逆映射的时间复杂度为 $O(N_\theta \times N_h)$，在普通计算机上可在秒级完成A4幅面的纸面图案生成。闭式解的另一个优势在于可以直接计算雅可比矩阵，从而精确分析映射的局部畸变特性，为参数优化提供梯度信息。

模型创新性地引入了双目视觉因素约束，将瞳距、人眼等效焦距、清晰视区和舒适视区等人因参数纳入设计框架。通过推导虚像深度与双目视差角的定量关系，建立了虚像平面一致性约束，为圆柱参数的选取提供了基于人因工程的设计准则。这一扩展使模型从纯几何光学层面提升到了考虑人眼感知特性的层面，更贴近实际观察体验。

多目标优化框架兼顾图像质量、畸变均匀性和边界约束三个维度，通过权重系数 $\lambda_1, \lambda_2, \lambda_3$ 的调节可灵活适应不同设计需求。当设计者追求镜面图案的高保真度时，可增大 $\lambda_1$；当设计者希望纸面图案的畸变更均匀（便于手绘或印刷）时，可增大 $\lambda_2$。这种灵活性使得模型能够适应从精密印刷到手工绘制的不同应用场景。

兼容性指标体系从频谱、对称性、容错度、主方向四个维度综合评价图案组合的可行性，具有较强的理论解释力和实践指导意义。该指标体系不仅可以判断给定图案组合的可行性等级，还可以为设计者提供改进方向——例如，当 $C_{\text{sym}}$ 较低时，建议选择对称性更强的图案；当 $C_{\text{freq}}$ 较低时，建议降低图案的空间频率复杂度。

\subsection{模型局限}

本文模型在以下方面存在局限性，需要在后续研究中加以改进。

远场平行光近似在观察距离较近时误差增大。当观察距离 $L$ 与圆柱半径 $R$ 的比值 $L/R < 20$ 时，平行光近似引入的角度误差超过3\%，对于近距离观察场景（如桌面摆件）需要引入透视投影模型，将入射光线从平行光推广为从观察者位置出发的发散光线。透视模型下逆映射不再具有闭式解，需要采用数值射线追踪方法。

模型仅考虑一次反射，忽略了多次反射和散射效应。在高反射率环境中，光线可能在圆柱面和纸面之间发生多次反射，产生"鬼影"效应。对于反射率低于90\%的普通镜面，二次反射的能量衰减至原始能量的1\%以下，可以忽略；但对于高反射率镜面（如镀银圆柱），多次反射可能对图案质量产生可观测的影响。

SSIM指标主要反映结构相似性，对颜色失真的敏感度有限。在彩色图案设计中，SSIM可能无法准确反映人眼对色彩差异的感知，需要引入感知颜色差异指标（如CIEDE2000）或多通道SSIM（MS-SSIM）作为补充评价指标。此外，SSIM对图像的全局亮度偏移不敏感，而圆柱反射可能引入系统性的亮度变化。

\subsection{模型推广}

本文模型可在以下方向进行推广，拓展圆柱镜反射艺术的应用空间。

将圆柱面替换为球面、抛物面、椭球面等其他曲面，推导相应的逆映射公式，可设计更丰富的反射艺术作品。球面镜的逆映射具有径向对称性，适合设计圆形图案；抛物面镜具有聚焦特性，可以实现局部放大效果。Luycho倒影杯的曲面设计即为此类推广的实际应用，其杯身曲面经过精心设计，使得碟子上的图案经杯身反射后呈现出丰富的视觉效果。

通过设计随观察角度变化而呈现不同镜面图案的纸面图案，可实现动态视觉效果。这需要在多个观察角度 $(\phi_1, \alpha_1), (\phi_2, \alpha_2), \ldots$ 下同时优化逆映射，是一个多目标、多约束的优化问题。动态反射艺术在互动展览和教育领域具有广阔的应用前景。

将纸面从平面推广到曲面（如球面、圆柱面），可设计更复杂的立体反射艺术作品。此时逆映射需要考虑纸面曲率对反射光线与纸面交点的影响，几何关系更为复杂，但基本的反射定律和映射框架仍然适用。三维纸面图案与曲面镜的组合可以创造出更加丰富的视觉体验，进一步拓展反射艺术的表现力。
